\section{Related Work}
\label{sec:related}

\paragraph{World Models.}
The concept of a world model as an internal simulator that predicts future states given actions is central to several frameworks for general intelligence~\citep{lecun2022path, hafner2023mastering, yang2026world, delgrange2026foundation}.
In visual and embodied domains, this idea has been instantiated through learned latent-space dynamics models.
IRIS~\citep{micheli2022transformers} first demonstrates that Transformers can serve as sample-efficient world models for Atari games.
DreamerV3~\citep{hafner2023mastering} and its successor Dreamer~4~\citep{hafner2025training} train agents entirely inside a learned world model, achieving strong performance across diverse control tasks without direct environment interaction during policy optimization.
Video-based world models such as Cosmos~\citep{ali2025world}, Genie~3~\citep{genie3}, and Vid2World~\citep{huang2025vid2world} scale this approach to high-resolution visual prediction, enabling interactive simulation of physical environments.
GameNGen~\citep{valevski2025diffusion} shows that diffusion models can simulate complex video games in real time.
V-JEPA~2~\citep{assran2025v} learns predictive representations from video without pixel-level reconstruction, demonstrating that self-supervised world models can support both understanding and planning.
UniSim~\citep{yang2023learning} learns a universal simulator of real-world interaction by orchestrating diverse visual datasets spanning objects, robotic actions, and navigation.
\citet{hou2026world} provide a comprehensive survey of world models for robot learning, covering policy coupling and video world models across manipulation, navigation, and autonomous driving.
More broadly, these approaches model how the environment evolves over time, often conditioned on actions, and use the learned dynamics for planning or policy learning.
\method extends this paradigm to text-based agent environments, where the ``observations'' are structured text such as API responses, terminal outputs, accessibility trees, and UI view hierarchies.

\paragraph{Language World Models.}
A growing body of work explores whether LLMs can serve as world models for text-based agent environments~\citep{li2026bridgingagentworldgaptext}.
\citet{li2025word} propose a three-level evaluation framework for LLM-based world models and find that fine-tuning substantially improves simulation fidelity, though gains depend on behavioral coverage and environment complexity; \citet{wang2024can} draw a complementary conclusion through systematic evaluation of LLMs as text-game simulators, showing that off-the-shelf LLMs remain unreliable world simulators.
\citet{richens2025general} prove theoretically that any sufficiently general agent must contain a world model, and \citet{cifuentes2026general} extend this result to partial observability and stochastic settings.
On the empirical side, several concurrent works train LLMs as environment simulators.
RLVR-World~\citep{wu2026rlvr} uses RL to train a unified world model that generates environment responses for agent training across text, web, and video domains, demonstrating that RL-trained simulators produce higher-fidelity predictions than SFT-only baselines.
In the web domain, WebDreamer~\citep{gu2024your} and WMA~\citep{chae2025web} pioneer model-based planning with LLM-simulated web environments, WebWorld~\citep{xiao2026webworld} is the first open-web world model trained at scale advancing both agent training and inference-time search, and WebATLAS~\citep{cheng2025webatlas} combines experience-driven memory with look-ahead action simulation.
\citet{shen2026world} augment web agents with world-model-based action correction, Imagine-then-Plan~\citep{liu2026imagine} uses world-model lookahead for adaptive agent planning, DynaWeb~\citep{ding2026dynaweb} applies model-based RL to web agents, and WebSynthesis~\citep{gao2025websynthesis} combines world-model-guided MCTS with trajectory synthesis.
Simia~\citep{li2025simulating} uses reasoning models as environment simulators for agent training on the $\tau$-bench task suite.
RWML~\citep{yu2026reinforcement} formulates world-model learning as an RL problem and demonstrates that RL-trained world-modeling capability improves the model's own agent performance.
DyMo~\citep{guo2025world} shows that language world modeling improves agent performance on interactive tasks, validating the benefit of next-state prediction as a training objective.
\citet{wang2025llms} show that LLMs can serve as scalable, general-purpose simulators for digital agent training.
\citet{zhang2025agent} propose learning from \emph{early experience}, where the agent uses its own interaction data with implicit world modeling as supervision, bypassing the need for external rewards.
On the training methodology front, VAGEN~\citep{wang2026vagen} reinforces world-model reasoning in VLM agents through dense state-prediction rewards, BehR~\citep{huang2026beyond} optimizes for behavior consistency with the real environment rather than surface-level state matching, \citet{ren2026aligning} align agentic world models via knowledgeable experience learning, and SWIRL~\citep{qiu2026self} self-improves world models from state-only sequences by alternating forward and inverse dynamics models without action labels.
SSRL~\citep{fan2025ssrl} and ZeroSearch~\citep{sun2025zerosearch} train LLMs via RL to simulate search engines, eliminating dependence on external APIs.
ECHO~\citep{shrivastava2026echo} shows that terminal agents can acquire world models by adding an auxiliary environment-prediction loss to standard RL, doubling pass rates on Terminal-Bench~2.0.
In the GUI domain, Code2World~\citep{zheng2026code2world} generates renderable code as a proxy for GUI state prediction, while CUWM~\citep{guan2026computer} builds a computer-use world model via two-stage factorization.
MobileDreamer~\citep{cao2026mobiledreamer} uses sketch-based GUI world models for mobile agent planning, \citet{koh2026generative} propose generative visual code world models for mobile environments, and UISim~\citep{xiang2025uisim} builds an interactive image-based UI simulator for dynamic mobile environments.
SWE-World~\citep{sun2026swe} proposes LLM-based environment simulation for software engineering agents, removing the dependence on Docker sandboxes, CWM~\citep{copet2025cwm} mid-trains a 32B open-weights LLM on observation-action trajectories from code interpreters and Docker environments to enable step-by-step simulation of code execution, and \citet{rahmani2026debugging} analyze and debug failure modes in code world models.
A complementary line uses LLMs to generate executable world models rather than serving as simulators directly: Code~WM~\citep{lehrach2025code} translates game rules into Python simulators supporting MCTS planning, Text2World~\citep{hu2025text2world} benchmarks LLMs on PDDL-based symbolic world model generation, and Agent2World~\citep{hu2025agent2world} proposes a multi-agent framework for the same task.
WorldLLM~\citep{levy2025worldllm} takes an orthogonal approach, enhancing world modeling through curiosity-driven theory-making without fine-tuning.
\citet{qian2026current} provide an important negative result: current agents fail to leverage world models as tools for foresight-based planning, lacking the mechanisms to strategically invoke simulation and integrate its predictions.
This finding motivates our native world-model approach, where environment modeling is the training objective from the CPT stage onward rather than bolted on after training.

\method differs from these works in two respects.
First, we develop a \emph{native} language world model as a foundation model for agentic environment simulation: it covers 7 domains within a single model, and is trained through a three-stage CPT$\to$SFT$\to$RL pipeline that starts from environment modeling as the pre-training objective, rather than fine-tuning a general-purpose LLM post hoc.
Second, we investigate how world modeling can improve general agents through two complementary paradigms: \emph{decoupling} the agent from the world model, where turn-level controllability (partial observation, difficulty modulation) enables Sim~RL that matches or exceeds Real~RL; and \emph{unifying} them into a single framework, where LWM pre-training serves as a warm-up that strengthens downstream agent performance.

\paragraph{Agent Foundation Models and Continual Pre-Training.}
AgentFounder~\citep{su2025scaling} proposes continual pre-training on agentic interaction data to build a 30B-parameter agent foundation model, demonstrating that domain-specific CPT improves downstream agent performance.
TRUSTEE~\citep{tang2026democratizingtoollearningenvironments} shows that even a free 8B LLM can fully simulate tool-use environments, enabling zero-shot agent training without any real environment access.
\citet{chen2025internalizing} internalize world models via self-play fine-tuning. In the robotics domain, DreamZero~\citep{ye2026world} shows that video-based world-action models trained on manipulation data can serve as zero-shot policies.
Our work shares the CPT-stage insight with AgentFounder but differs in scope (seven domains vs.\ a smaller set) and in the subsequent SFT and RL stages that refine simulation fidelity beyond what CPT alone achieves.

\paragraph{Synthetic Environment Generation.}
A parallel line of work constructs synthetic environments programmatically rather than learning a neural simulator.
AWM (Agent World Model)~\citep{wang2026agent} generates code-based environments for agentic RL, and Agent-World~\citep{dong2026agent} scales real-world environment synthesis to nearly 2,000 environments with self-evolving training.
RLAnything~\citep{wang2026rlanything} jointly forges environment, policy, and reward model in a fully dynamic RL system, and GenEnv~\citep{guo2025genenv} co-evolves agents and environment simulators with difficulty-aligned curricula.
ASTRA~\citep{tian2026astra} automates the synthesis of agentic trajectories and reinforcement arenas from tool-call graphs.
Domain-specific generators include InfiniteWeb~\citep{zhang2026infiniteweb}, WebGym~\citep{bai2026webgym}, Weblica~\citep{kar2026weblicascalablereproducibletraining}, AutoWebWorld~\citep{wu2026autowebworld}, WebFactory~\citep{fan2026webfactory}, and \citet{chae2026safe} for web environments; GUI-Genesis~\citep{cao2026gui}, EvoCUA~\citep{xue2026evocua}, and EE-MCP~\citep{he2026ee} for GUI and mobile environments; TermiGen~\citep{zhu2026termigen} and Endless Terminals~\citep{gandhi2026endless} for terminal environments; and SWE-Universe~\citep{chen2026swe}, daVinci-Env~\citep{fu2026davinci}, \citet{zhao2026immersion}, SandMLE~\citep{zhou2026synthetic}, and ClawEnvKit~\citep{li2026clawenvkit} for software engineering and general agent platforms.
ScaleEnv~\citep{tu2026scaleenv}, AutoForge~\citep{cai2025autoforge}, EnvScaler~\citep{song2026envscaler}, SCALER~\citep{xu2026scaler}, and \citet{fang2025towards} propose general frameworks for scaling environment synthesis.
These code-driven approaches offer deterministic execution and verifiable rewards, but are inherently limited to domains where environments can be programmatically specified.
Our LWM-based simulation is complementary: it trades determinism for generality, covering domains (e.g., search engines, real-world MCP servers) where code-based synthesis is impractical.
\citet{huang2025environment} survey environment scaling methods for interactive agentic experience collection, and \citet{chu2026agentic} provide a broader treatment of agentic world modeling covering foundations, capabilities, and emerging directions.
